% ============================================================================
% Section 1 — Introduction + Related Work (mock-up)
% References for the \cite keys live in sections/99_references.tex
% ============================================================================
% PLANNING NOTES — key arguments the author must develop
% ============================================================================
% A1. FRAMING: The scale-vehicle hybrid (fuel cell + battery on a shared DC bus)
%     is a *miniature DC microgrid*. This lets us import the whole droop /
%     current-sharing literature wholesale, and is the justification for
%     submitting to Energies rather than a vehicles venue. Say this explicitly
%     in the first paragraph of the Introduction.
%
% A2. THE GAP (this is the paper's actual novelty — hammer it):
%     Nearly all droop-sharing literature assumes a *custom-designed* converter
%     whose inner current/voltage loop is authored by the researcher (DSP/FPGA
%     PWM generation, full state access). Practical builders of small hybrid
%     systems instead use monolithic, internally-compensated boost ICs
%     (TPS61288 class) where the control loop is inaccessible — the only exposed
%     node is the FB divider. The literature is essentially silent on how to get
%     *commanded, continuously variable* power sharing out of such parts.
%     => Claim: droop realized by MDAC perturbation of the FB node of an
%     off-the-shelf boost IC. Cost/effort of a custom converter avoided;
%     controllability of a centralized scheme retained.
%
% A3. WHY NOT JUST DO IT IN SOFTWARE / COMMUNICATION:
%     Argue that the FC and battery branches are physically separate sources;
%     a comms-based master-slave scheme adds a link, a single point of
%     failure, and latency inside the fast current loop. Droop keeps the fast
%     sharing action analog and local; the slow outer loop only *biases* it.
%     This is the hierarchical-control argument (Guerrero et al.).
%
% A4. WHY ROBUST CONTROL RATHER THAN A PI:
%     The plant seen by the outer share loop is *not* well known: it is the
%     product of (i) IC-internal compensation the vendor only publishes
%     approximately, (ii) INA sense-gain and divider tolerance, (iii) a
%     bus-voltage-dependent and load-current-dependent operating point, and
%     (iv) a droop-injection network with its own attenuation. State the
%     parameter uncertainty box quantitatively (cite the plant model section)
%     and argue this is exactly the setting an H-inf / Youla design addresses,
%     whereas a hand-tuned PI is only valid at one corner.
%     -> Back with the concrete result: worst-corner ||S||_inf 1.87 vs 26.9,
%     all plant corners stable, delay margin ~11.7 ms. MOVE THESE NUMBERS
%     to the Results section but *promise* them in the Introduction.
%
% A5. HONEST LIMITATIONS to pre-empt reviewers (state in Intro, revisit in
%     Discussion): droop trades bus-voltage regulation for sharing accuracy;
%     the MDAC path adds an extra tolerance stack; the approach is bounded by
%     the boost IC's own stability (RHP zero, compensation) so it does not
%     generalize to arbitrary converters without a per-part plant model.
%
% A6. CONTRIBUTION LIST (end of Introduction, 4 bullets — keep parallel):
%     (1) an FB-injection droop mechanism for internally-compensated boost ICs;
%     (2) a small-signal plant model of the two-converter droop-sharing loop
%         including injection attenuation and sense chain, validated against a
%         full-order model;
%     (3) a robust (H-inf/Youla) outer share controller with a documented
%         synthesis and gate-check procedure;
%     (4) experimental validation on a scale fuel-cell/battery vehicle.
%
% A7. TODO(author): decide whether to cite the scale-vehicle / RC-testbed
%     literature at all. If a suitable "hardware-in-the-loop scale testbed"
%     reference exists, add it to justify scale validation as scientifically
%     meaningful. Currently no such citation is included.
%
% A8. TODO(verify): every \bibitem in 99_references.tex was checked to exist
%     and carries a DOI or stable URL, but several carry % TODO(verify) marks
%     for author list / volume — re-check against Crossref before submission;
%     MDPI copy-edits references against Crossref.
% ============================================================================

\section{Introduction}
\label{sec:intro}

% NOTE: Paragraph 1 should establish the application context and assert the
% microgrid framing (A1). Keep it short; do not oversell hydrogen.
Hybridizing a fuel cell with a battery is the standard remedy for the fuel
cell's poor transient response and its intolerance of load steps and reverse
current. In such a powertrain both sources feed a common DC link through
dedicated converters, and the question of how instantaneous load current is
apportioned between them becomes a control problem in its own right. Structurally
this is the same problem studied under the heading of DC microgrids, where
multiple dissimilar sources share a bus through parallel-connected
converters~\cite{dragicevic2016part1,dragicevic2016part2,gao2018primary}.

% NOTE: Paragraph 2 introduces droop as the incumbent decentralized answer and
% states its two well-known weaknesses (voltage deviation, sharing error under
% line/parameter mismatch). This sets up why an outer loop is needed at all.
The conventional answer is droop control: each converter deliberately reduces
its output voltage in proportion to its output current, so that a virtual output
resistance sets the sharing ratio without any communication between
units~\cite{guerrero2011hierarchical,tah2016adaptive}. Droop is attractive because
it is local, fast and inherently redundant, but it is well documented that it
trades bus-voltage regulation against sharing accuracy, and that mismatched line
and converter parameters corrupt the achieved split~\cite{anand2013distributed,lu2014adaptive}.
A large body of work therefore adds a secondary or distributed layer that trims
the droop coefficients or the voltage
references~\cite{nasirian2015cooperative,tucci2018plugandplay,kumar2020distributed}.

% NOTE: Paragraph 3 is the pivot — the hardware-realization gap (A2). This is
% the paragraph the whole paper hangs on. Author should be blunt: existing
% droop work presumes a converter you designed yourself.
Almost without exception, that literature presumes a converter whose inner loop
the designer owns: modulation is generated in a DSP or FPGA, and the droop term
is simply added to a software voltage reference. Practitioners building small
hybrid systems rarely enjoy that freedom. Power stages are increasingly bought as
monolithic, internally-compensated regulator ICs, in which the control loop is
sealed and the only user-accessible node is the feedback divider. For such parts
the standard droop recipe has no point of entry, and designers fall back on
fixed resistive droop, on a hard master--slave split, or on abandoning
controllable sharing altogether.

% NOTE: Paragraph 4 states the proposed mechanism concretely. Do NOT get into
% circuit values here — forward-reference the hardware section.
This paper shows that controllable droop can nevertheless be obtained from such
a device. A multiplying digital-to-analogue converter (MDAC) injects a
current-proportional perturbation into the feedback node of each off-the-shelf
boost regulator, synthesizing a commanded, continuously variable output
resistance in analogue hardware. The fast sharing action remains local to each
converter; a slow digital outer loop merely commands the droop ratio, and hence
the steady-state power split, between the fuel-cell and battery branches.

% NOTE: Paragraph 5 motivates the robust outer controller (A4). Emphasize
% *why the plant is uncertain*, not just that robust control is fashionable.
Because the resulting plant is assembled from an IC whose internal compensation
is only approximately published, a sense amplifier and divider network with
finite tolerance, and an operating point that moves with bus voltage and load,
its parameters are known only within a box rather than exactly. This is precisely
the setting for robust synthesis, and $\mathcal{H}_\infty$ and related methods
have already been applied to microgrid voltage and sharing
loops~\cite{sedhom2019hinf,hosseinzadeh2024robust,mirjafari2022robust}. We
accordingly design the outer share controller by $\mathcal{H}_\infty$/Youla
synthesis over the parameter corners rather than by tuning a PI at nominal.

% NOTE: Paragraph 6 = contributions. Use A6. Keep the four items parallel.
The contributions of this work are: (i) a feedback-injection droop mechanism
that gives commanded power sharing from internally-compensated, off-the-shelf
boost regulators; (ii) a small-signal model of the two-converter droop-sharing
loop that accounts for the injection attenuation and the current-sense chain,
validated against an independently constructed full-order model; (iii) a robust
outer share controller with a documented synthesis and verification procedure;
and (iv) experimental validation on a scale fuel-cell/battery vehicle.

% NOTE: Paragraph 7 = roadmap. Author fills in once section numbering is final.
The remainder of the paper is organized as follows. Section~\ref{sec:related}
reviews power-sharing methods and positions the proposed approach.
Section~\ref{sec:droop} develops the droop design and gain selection;
Section~\ref{sec:robust} the robust share controller; Section~\ref{sec:board-design}
the board and power-path hardware; Section~\ref{sec:bringup} the bring-up
diagnosis; Section~\ref{sec:mimo} the MIMO outlook.

\section{Related Work}
\label{sec:related}

\subsection{Droop Control and Its Refinements}

% NOTE: Establish the canonical V-I droop / virtual resistance formulation and
% its known limits. Cite the two survey papers so a reviewer sees breadth.
Conventional voltage--current droop emulates a series resistance $R_d$ at each
converter output, so that $v_o = V^\ast - R_d i_o$. Sharing accuracy improves
with larger $R_d$ while bus regulation degrades, a tension that is the
starting point of most subsequent work~\cite{dragicevic2016part1,gao2018primary}.
Virtual-resistance implementations move this behaviour into the controller rather
than dissipating it, and adaptive schemes vary $R_d$ online according to
converter rating, state of charge or measured sharing
error~\cite{tah2016adaptive,lu2014adaptive,khorsandi2024adaptive}.

% NOTE: Second paragraph should cover the secondary/distributed correction
% layer, and make the point that these schemes reintroduce communication —
% which is the wedge for our comparison table.
Because purely local droop cannot remove the steady-state voltage deviation, a
secondary layer is usually added. Distributed and consensus-based schemes
exchange voltage and current estimates over a sparse communication graph to
restore the bus and equalize per-unit currents
simultaneously~\cite{nasirian2015cooperative,anand2013distributed,kumar2020distributed},
and plug-and-play designs provide stability guarantees that are independent of
the network topology~\cite{tucci2018plugandplay}. These methods achieve excellent
accuracy at the cost of a communication link and its associated latency and
failure modes.

\subsection{Alternatives to Droop}

% NOTE: This subsection must be genuinely fair to the alternatives — reviewers
% punish strawmanning. Cover master-slave, average current sharing, DC bus
% signaling, MPC. Each gets one or two sentences plus a citation.
Active current-sharing methods predate the microgrid framing. Master--slave and
average-current-sharing architectures for paralleled supply modules were
classified early and remain the industrial default where the modules are
identical and co-located~\cite{luo1999classification,panov2008stability}. They
give near-exact sharing but require a dedicated share bus and degrade poorly
when the master fails.

DC bus signaling dispenses with an explicit link by encoding mode transitions in
the bus voltage itself, so that each unit infers the system state
locally~\cite{schonberger2006dbs,panda2021dbs}. The method is elegant for
coarse mode arbitration but does not deliver a continuously commandable split.

Optimization-based supervision, particularly model predictive control, has been
applied extensively to fuel-cell hybrid powertrains to schedule the source split
against fuel economy and degradation
objectives~\cite{vermillion2011mpc,hu2021mpcreview,zhou2022mpc}. Such strategies
operate above the converter loops: they decide \emph{what} the split should be, and
still require an actuation mechanism to enforce it. The present work is
complementary rather than competing, and in fact supplies exactly that
mechanism.

\subsection{Power Sharing in Fuel-Cell Hybrid Vehicles}

% NOTE: Ground the work in the FC/battery application; make the frequency-
% separation argument (FC slow, battery fast) since it explains why the share
% setpoint is slow-moving and why a slow robust outer loop suffices.
In fuel-cell hybrids the split is not arbitrary: the stack must be protected
from fast transients, which are absorbed by the battery or a supercapacitor,
while the stack carries the slow-moving mean
demand~\cite{thounthong2007control,garcia2016statemachine,li2022dcbusregulation}.
This frequency separation is usually enforced by filtering the demand signal or
by a rule-based supervisor. Droop achieves the same separation as a structural
property when the branches are given different output impedances, which is one
reason droop-based energy management has been reported for fuel-cell
tramways and hybrid buses~\cite{garcia2016statemachine,li2022dcbusregulation}.

\subsection{Robust Control of Sharing Loops}

% NOTE: Short subsection. Point is: robust methods exist for microgrid loops but
% are applied to custom converters with accessible states; nobody applies them
% to a droop loop whose actuator is an MDAC on a commercial IC's FB pin.
Robust synthesis has been used to certify microgrid loops against parameter
drift, constant-power-load negative impedance and communication
delay~\cite{sedhom2019hinf,hosseinzadeh2024robust,mirjafari2022robust}. These
results confirm the value of a worst-case formulation, but they are invariably
applied to converters whose internal dynamics are designed alongside the
controller. Applying the same discipline to a loop whose actuator is an MDAC
perturbing the feedback pin of a commercial regulator, and whose plant gain
therefore inherits vendor tolerances, appears not to have been reported.

\subsection{Positioning}

% NOTE: Refer to the table explicitly and state the one-line thesis of the paper.
Table~\ref{tab:comparison} summarizes the trade-offs. The proposed approach
occupies a cell that the existing methods do not: it retains the commandability
of a centralized scheme and the robustness of a droop scheme, while requiring
neither a custom power stage nor a communication link inside the fast loop.

% ---------------------------------------------------------------------------
% Comparison table
% ---------------------------------------------------------------------------
\begin{table}[H]
\caption{Comparison of DC power-sharing methods against the criteria relevant to
a low-cost hybrid powertrain. Ratings are qualitative:
\textbullet\textbullet\textbullet{} = favourable,
\textbullet{} = unfavourable.}
\label{tab:comparison}
\newcommand{\rgood}{\textbullet\textbullet\textbullet}
\newcommand{\rmid}{\textbullet\textbullet}
\newcommand{\rbad}{\textbullet}
\begin{adjustwidth}{-\extralength}{0cm}
\begin{tabularx}{\fulllength}{lCCCCCl}
\toprule
\textbf{Method} & \textbf{Hardware} & \textbf{Comms} & \textbf{Split} &
\textbf{Robustness to} & \textbf{Off-the-shelf} & \textbf{Refs} \\
 & \textbf{cost} & \textbf{required} & \textbf{controllability} &
\textbf{converter tol.} & \textbf{ICs usable} & \\
\midrule
Fixed resistive droop        & \rgood & none    & \rbad  & \rbad  & \rgood & \cite{dragicevic2016part1,gao2018primary} \\
Virtual-resistance droop     & \rmid  & none    & \rmid  & \rmid  & \rbad  & \cite{tah2016adaptive,lu2014adaptive} \\
Adaptive / consensus droop   & \rmid  & sparse  & \rgood & \rgood & \rbad  & \cite{nasirian2015cooperative,anand2013distributed,kumar2020distributed} \\
Master--slave current share  & \rbad  & share bus & \rgood & \rmid & \rbad  & \cite{luo1999classification} \\
Average current sharing      & \rbad  & share bus & \rmid & \rmid  & \rbad  & \cite{luo1999classification,panov2008stability} \\
DC bus signaling             & \rgood & none    & \rbad  & \rmid  & \rmid  & \cite{schonberger2006dbs,panda2021dbs} \\
Supervisory MPC / EMS        & \rmid  & to supervisor & \rgood & n/a $^{a}$ & n/a $^{a}$ & \cite{vermillion2011mpc,hu2021mpcreview,zhou2022mpc} \\
\midrule
\textbf{This work:} MDAC feedback-injection droop
                             & \rgood & none $^{b}$ & \rgood & \rgood $^{c}$ & \rgood & --- \\
\bottomrule
\end{tabularx}
\end{adjustwidth}
\noindent{\footnotesize
$^{a}$ Supervisory strategies presuppose an underlying actuation mechanism and
are therefore complementary to, not substitutes for, the rows above.
$^{b}$ No communication inside the fast sharing loop; a low-rate link carries
only the share setpoint.
$^{c}$ By construction of the robust outer controller over the parameter
corners; see Section~\ref{sec:robust}.}
\end{table}

% NOTE(author): the table uses MDPI's adjustwidth/tabularx macros from the
% Energies class (\fulllength, \extralength, column type C). The agent draft
% used threeparttable-style \tablenotes, which the MDPI class does NOT load —
% replaced here with plain footnote markers. Verify column type C compiles;
% if not, define \newcolumntype{C}{>{\centering\arraybackslash}X}.
